%%% For single-sided printing
%\documentclass[11pt,a4paper]{report}
%\usepackage[top=25mm,bottom=25mm,right=25mm,left=30mm,head=12.5mm,foot=12.5mm]{geometry}
%\let\openright=\clearpage

%%% For two-sided printing:
\documentclass[11pt,a4paper,twoside,openright]{report}
\usepackage[top=25mm,bottom=25mm,right=25mm,left=30mm,head=12.5mm,foot=12.5mm]{geometry}
\let\openright=\cleardoublepage

%%% References
\usepackage{hyperref}
\usepackage{nameref}

%% Description reference labels
\makeatletter
\def\namedlabel#1#2{\begingroup
    #2%
    \def\@currentlabel{#2}%
    \phantomsection\label{#1}\endgroup
}
\makeatother

%%% Charts
\usepackage{pgf-pie}
\usepackage{pgfplots}

%%% Footnotes
\usepackage{footnote}
\makesavenoteenv{tabular}

%%% Definition of useful macros
%%% Tento soubor obsahuje definice různých užitečných maker a prostředí %%%
%%% Další makra připisujte sem, ať nepřekáží v ostatních souborech.     %%%

\usepackage[a-2u]{pdfx}     % výsledné PDF bude ve standardu PDF/A-2u

\usepackage{ifpdf}
\usepackage{ifxetex}
\usepackage{ifluatex}

%%% Nastavení pro použití samostatné bibliografické databáze.
\usepackage[
   backend=biber
%  ,style=iso-authoryear
  ,style=iso-numeric
  ,sortlocale=cs_CZ
  ,alldates=iso
  ,bibencoding=UTF8
  %,block=ragged
]{biblatex}
\let\cite\parencite
\bibliography{literatura}

%% Přepneme na českou sazbu, fonty Latin Modern a kódování češtiny
\ifthenelse{\boolean{xetex}\OR\boolean{luatex}}
   { % use fontspec and OpenType fonts with utf8 engines
			\usepackage[english,slovak,czech]{babel}
			\usepackage[autostyle,english=british,czech=quotes]{csquotes}
			\usepackage{fontspec}
			\defaultfontfeatures{Ligatures=TeX,Scale=MatchLowercase}
   }
   {
			\usepackage[english,slovak,czech]{babel}
			\usepackage{lmodern}
			\usepackage[T1]{fontenc}
			\usepackage{textcomp}
			\usepackage[utf8]{inputenc}
			\usepackage[autostyle,english=british,czech=quotes]{csquotes}
	 }
\ifluatex
\makeatletter
\let\pdfstrcmp\pdf@strcmp
\makeatother
\fi

%%% Další užitečné balíčky (jsou součástí běžných distribucí LaTeXu)
\usepackage{amsmath}        % rozšíření pro sazbu matematiky
\usepackage{amsfonts}       % matematické fonty
\usepackage{amssymb}        % symboly
\usepackage{amsthm}         % sazba vět, definic apod.
\usepackage{bm}             % tučné symboly (příkaz \bm)
\usepackage{graphicx}       % vkládání obrázků
\usepackage{listings}       % vylepšené prostředí pro strojové písmo
\usepackage{fancyhdr}       % prostředí pohodlnější nastavení hlavy a paty stránek
\usepackage{icomma}         % inteligetní čárka v matematickém módu
\usepackage{dcolumn}        % lepší zarovnání sloupců v tabulkách
\usepackage{booktabs}       % lepší vodorovné linky v tabulkách
\makeatletter
\@ifpackageloaded{xcolor}{
   \@ifpackagewith{xcolor}{usenames}{}{\PassOptionsToPackage{usenames}{xcolor}}
  }{\usepackage[usenames]{xcolor}} % barevná sazba
\makeatother
\usepackage{multicol}       % práce s více sloupci na stránce
\usepackage{caption}
\usepackage{enumitem}
\setlist[itemize]{noitemsep, topsep=0pt, partopsep=0pt}
\setlist[enumerate]{noitemsep, topsep=0pt, partopsep=0pt}
\setlist[description]{noitemsep, topsep=0pt, partopsep=0pt}

\usepackage{tocloft}
\setlength\cftparskip{0pt}
\setlength\cftbeforechapskip{1.5ex}
\setlength\cftfigindent{0pt}
\setlength\cfttabindent{0pt}
\setlength\cftbeforeloftitleskip{0pt}
\setlength\cftbeforelottitleskip{0pt}
\setlength\cftbeforetoctitleskip{0pt}
\renewcommand{\cftlottitlefont}{\Huge\bfseries\sffamily}
\renewcommand{\cftloftitlefont}{\Huge\bfseries\sffamily}
\renewcommand{\cfttoctitlefont}{\Huge\bfseries\sffamily}

% vyznaceni odstavcu
\parindent=0pt
\parskip=11pt

% zakaz vdov a sirotku - jednoradkovych pocatku ci koncu odstavcu na prechodu mezi strankami
\clubpenalty=1000
\widowpenalty=1000
\displaywidowpenalty=1000

% nastaveni radkovani
\renewcommand{\baselinestretch}{1.20}

% nastaveni pro nadpisy - tucne a bezpatkove
\usepackage{sectsty}    
\allsectionsfont{\sffamily}

% nastavení hlavy a paty stránek
\fancyhf{}
\fancyhead[RO,LE]{\rightmark}
\fancyfoot[RO,LE]{\thepage}
\renewcommand{\footrulewidth}{.5pt}
\fancypagestyle{plain}{%
\fancyhf{} % clear all header and footer fields
\fancyfoot[RO,LE]{\thepage}
\renewcommand{\headrulewidth}{0pt}
\renewcommand{\footrulewidth}{0.5pt}}

% Tato makra přesvědčují mírně ošklivým trikem LaTeX, aby hlavičky kapitol
% sázel příčetněji a nevynechával nad nimi spoustu místa. Směle ignorujte.
\makeatletter
\def\@makechapterhead#1{
  {\parindent \z@ \raggedright \sffamily
   \Huge\bfseries \thechapter. #1
   \par\nobreak
   \vskip 20\p@
}}
\def\@makeschapterhead#1{
  {\parindent \z@ \raggedright \sffamily
   \Huge\bfseries #1
   \par\nobreak
   \vskip 20\p@
}}
\makeatother

% Trochu volnější nastavení dělení slov, než je default.
\lefthyphenmin=2
\righthyphenmin=2

% Zapne černé "slimáky" na koncích řádků, které přetekly, abychom si
% jich lépe všimli.
\overfullrule=1mm

%% Balíček hyperref, kterým jdou vyrábět klikací odkazy v PDF,
%% ale hlavně ho používáme k uložení metadat do PDF (včetně obsahu).
%% Většinu nastavítek přednastaví balíček pdfx.
\hypersetup{unicode}
\hypersetup{breaklinks=true}
\hypersetup{hidelinks}

%%% Prostředí pro sazbu kódu, případně vstupu/výstupu počítačových
%%% programů. (Vyžaduje balíček listings -- fancy verbatim.)
\lstnewenvironment{code}{\lstset{basicstyle=\small, frame=single}}{}





%%% DEFINICE ZÁKLADNÍCH PROMĚNNÝCH
\def\TypPrace{BP}               % bakalářská práce/bachelor thesis
%\def\TypPrace{DP}              % diplomová práce/master thesis
%\def\Jazyk{cze}                % čeština/czech
%\def\Jazyk{slo}                % slovenština/slovak
\def\Jazyk{eng}                 % angličtina/english

%%% Title of the thesis in the language used in the text (exact according to assignment)
\def\NazevPrace{Development of a documentation generating application for .NET libraries}

%%% Tento soubor obsahuje definice různých užitečných maker a prostředí %%%
%%% Další makra připisujte sem, ať nepřekáží v ostatních souborech.     %%%

\usepackage[a-2u]{pdfx}     % výsledné PDF bude ve standardu PDF/A-2u

\usepackage{ifpdf}
\usepackage{ifxetex}
\usepackage{ifluatex}

%%% Nastavení pro použití samostatné bibliografické databáze.
\usepackage[
   backend=biber
%  ,style=iso-authoryear
  ,style=iso-numeric
  ,sortlocale=cs_CZ
  ,alldates=iso
  ,bibencoding=UTF8
  %,block=ragged
]{biblatex}
\let\cite\parencite
\bibliography{literatura}

%% Přepneme na českou sazbu, fonty Latin Modern a kódování češtiny
\ifthenelse{\boolean{xetex}\OR\boolean{luatex}}
   { % use fontspec and OpenType fonts with utf8 engines
			\usepackage[english,slovak,czech]{babel}
			\usepackage[autostyle,english=british,czech=quotes]{csquotes}
			\usepackage{fontspec}
			\defaultfontfeatures{Ligatures=TeX,Scale=MatchLowercase}
   }
   {
			\usepackage[english,slovak,czech]{babel}
			\usepackage{lmodern}
			\usepackage[T1]{fontenc}
			\usepackage{textcomp}
			\usepackage[utf8]{inputenc}
			\usepackage[autostyle,english=british,czech=quotes]{csquotes}
	 }
\ifluatex
\makeatletter
\let\pdfstrcmp\pdf@strcmp
\makeatother
\fi

%%% Další užitečné balíčky (jsou součástí běžných distribucí LaTeXu)
\usepackage{amsmath}        % rozšíření pro sazbu matematiky
\usepackage{amsfonts}       % matematické fonty
\usepackage{amssymb}        % symboly
\usepackage{amsthm}         % sazba vět, definic apod.
\usepackage{bm}             % tučné symboly (příkaz \bm)
\usepackage{graphicx}       % vkládání obrázků
\usepackage{listings}       % vylepšené prostředí pro strojové písmo
\usepackage{fancyhdr}       % prostředí pohodlnější nastavení hlavy a paty stránek
\usepackage{icomma}         % inteligetní čárka v matematickém módu
\usepackage{dcolumn}        % lepší zarovnání sloupců v tabulkách
\usepackage{booktabs}       % lepší vodorovné linky v tabulkách
\makeatletter
\@ifpackageloaded{xcolor}{
   \@ifpackagewith{xcolor}{usenames}{}{\PassOptionsToPackage{usenames}{xcolor}}
  }{\usepackage[usenames]{xcolor}} % barevná sazba
\makeatother
\usepackage{multicol}       % práce s více sloupci na stránce
\usepackage{caption}
\usepackage{enumitem}
\setlist[itemize]{noitemsep, topsep=0pt, partopsep=0pt}
\setlist[enumerate]{noitemsep, topsep=0pt, partopsep=0pt}
\setlist[description]{noitemsep, topsep=0pt, partopsep=0pt}

\usepackage{tocloft}
\setlength\cftparskip{0pt}
\setlength\cftbeforechapskip{1.5ex}
\setlength\cftfigindent{0pt}
\setlength\cfttabindent{0pt}
\setlength\cftbeforeloftitleskip{0pt}
\setlength\cftbeforelottitleskip{0pt}
\setlength\cftbeforetoctitleskip{0pt}
\renewcommand{\cftlottitlefont}{\Huge\bfseries\sffamily}
\renewcommand{\cftloftitlefont}{\Huge\bfseries\sffamily}
\renewcommand{\cfttoctitlefont}{\Huge\bfseries\sffamily}

% vyznaceni odstavcu
\parindent=0pt
\parskip=11pt

% zakaz vdov a sirotku - jednoradkovych pocatku ci koncu odstavcu na prechodu mezi strankami
\clubpenalty=1000
\widowpenalty=1000
\displaywidowpenalty=1000

% nastaveni radkovani
\renewcommand{\baselinestretch}{1.20}

% nastaveni pro nadpisy - tucne a bezpatkove
\usepackage{sectsty}    
\allsectionsfont{\sffamily}

% nastavení hlavy a paty stránek
\fancyhf{}
\fancyhead[RO,LE]{\rightmark}
\fancyfoot[RO,LE]{\thepage}
\renewcommand{\footrulewidth}{.5pt}
\fancypagestyle{plain}{%
\fancyhf{} % clear all header and footer fields
\fancyfoot[RO,LE]{\thepage}
\renewcommand{\headrulewidth}{0pt}
\renewcommand{\footrulewidth}{0.5pt}}

% Tato makra přesvědčují mírně ošklivým trikem LaTeX, aby hlavičky kapitol
% sázel příčetněji a nevynechával nad nimi spoustu místa. Směle ignorujte.
\makeatletter
\def\@makechapterhead#1{
  {\parindent \z@ \raggedright \sffamily
   \Huge\bfseries \thechapter. #1
   \par\nobreak
   \vskip 20\p@
}}
\def\@makeschapterhead#1{
  {\parindent \z@ \raggedright \sffamily
   \Huge\bfseries #1
   \par\nobreak
   \vskip 20\p@
}}
\makeatother

% Trochu volnější nastavení dělení slov, než je default.
\lefthyphenmin=2
\righthyphenmin=2

% Zapne černé "slimáky" na koncích řádků, které přetekly, abychom si
% jich lépe všimli.
\overfullrule=1mm

%% Balíček hyperref, kterým jdou vyrábět klikací odkazy v PDF,
%% ale hlavně ho používáme k uložení metadat do PDF (včetně obsahu).
%% Většinu nastavítek přednastaví balíček pdfx.
\hypersetup{unicode}
\hypersetup{breaklinks=true}
\hypersetup{hidelinks}

%%% Prostředí pro sazbu kódu, případně vstupu/výstupu počítačových
%%% programů. (Vyžaduje balíček listings -- fancy verbatim.)
\lstnewenvironment{code}{\lstset{basicstyle=\small, frame=single}}{}





%%% Author's name - Firstname and Lastname
\def\AutorPrace{Denis Akopyan}

%%% Year of submission and month (verbally) - month YYYY
\def\DatumOdevzdani{měsíc 2022}

%%% Supervisor: First name and surname with titles
\def\Vedouci{Ing. Jarmila Pavlíčková, Ph.D.}

%%% Consultant: First name and surname with titles
\def\Konzultant{}

%%% Study program
\def\StudijniProgram{Applied Informatics}

%%% Study program - specialization
\def\Specializace{}

%%% Field of study
\def\StudijniObor{Applied Informatics}

%%% Optional thanks (the supervisor, the consultant, the borrower of software, literature, etc)
\def\Podekovani{%
Poděkování.
}

%%% Abstrakt (doporučený rozsah cca 150-250 slov; nejedná se o zadání práce)
\def\Abstrakt{%
Zaměřením dané bakalářské práce je výzkum a vývoj aplikace pro generování dokumentace .NET knihoven. Hlavním cílem bylo vytvoření nástroje pro .NET vývojáře, jež chtějí extrahovat dokumentaci zdrojového kódu a exportovat ji do požadovaného výstupu, a zároveň zachovat aplikaci rozšiřitelnou, multiplatformní a uživatelsky přístupnou.\

Teoretická část práce se bude zabývát plánováním vývoje a očekáváným výsledkem projektu.

Empirická část popíše vývoj samotný a zaměřuje se na tvorbu oddělené projektové struktury.

Očekáváným výsledkem bakalářské práce byde funkční, multiplatformní a open-source aplikace, která umožňuje generování dokumentace z .NET knihoven, a zároveň je rozšiřitelná prostřednictvím modulů.
}
\def\AbstraktEN{%
This thesis aims to research and develop a documentation-generating application for .NET libraries.
The main objective was to create a tool for .NET developers that want to extract source code documentation and export it into their desired format while keeping the application extensible, cross-platform, and easy to use.

The theoretical part of the thesis shall cover the development planning and the expected result of the
project.

The empirical part of the thesis will illustrate the exact process of development. Designing a decoupled
project structure had the most attention dedicated to it.

The expected result of this thesis project is a working cross-platform and open-source application that can
generate documentation from .NET libraries and is extensible via plugins.
}

%%% 3 až 5 klíčových slov (doporučeno)
\def\KlicovaSlova{vývoj, tvorba, aplikace, generování, dokumentace, dotnet, .NET, knihoven}
\def\KlicovaSlovaEN{development, application, generating, documentation, dotnet, .NET, libraries}

%%% Kody podle klasifikace JEL
\def\JEL{C88, O30}

%%% Fix height
\raggedbottom

%%% Title page and various mandatory information pages
\begin{document}
%%% Titulní strana práce a další povinné informační strany

%%% Titulní strana práce

\pagestyle{empty}
\hypersetup{pageanchor=false}

\begin{center}
\Huge\sffamily
\VSE\\
\FIS

\vspace{\stretch{1}}

\includegraphics[width=.3\textwidth]{img/logo-FIS}

\vspace{\stretch{2}}

\bfseries\NazevPrace

\vspace{8mm}
\mdseries\TypPraceText

\vspace{8mm}
\large
\begin{tabular}{rl}
\StudijniProgramText: & \StudijniProgram \\
\ifthenelse{\equal{\Specializace}{}}{%
	% empty value
	}{
	\rule{0pt}{6mm}%
	\SpecializaceText: & \Specializace \\
}
\ifthenelse{\equal{\StudijniObor}{}}{%
	% empty value
	}{
	\rule{0pt}{6mm}%
	\StudijniOborText: & \StudijniObor \\
}
\end{tabular}

\vspace{\stretch{8}}

\begin{tabular}{rl}
\AutorText: & \AutorPrace \\
\noalign{\vspace{2mm}}
\VedouciText: & \Vedouci \\
\ifthenelse{\equal{\Konzultant}{}}{%
	% empty value
	}{
	\rule{0pt}{6mm}%
	\KonzultantText: & \Konzultant \\
}
\end{tabular}

\vspace{8mm}
\Praha, \DatumOdevzdani
\end{center}


%%% Poděkování
\hypersetup{pageanchor=true}
\cleardoublepage
\pagestyle{plain}
\openright
\vspace*{\fill}
\section*{\PodekovaniText}
\noindent
\Podekovani
\vspace{1cm}


%%% Povinná informační strana bakalářské práce
\openright
\section*{Abstrakt}
\noindent
\Abstrakt
\subsection*{Klíčová slova}
\noindent
\KlicovaSlova
\subsection*{JEL klasifikace}
\noindent
\JEL

\newpage

\section*{Abstract}
\noindent
\AbstraktEN
\subsection*{Keywords}
\noindent
\KlicovaSlovaEN
\subsection*{JEL classification}
\noindent
\JEL

\openright


%%% A page with automatically generated content of the thesis
\setcounter{tocdepth}{2}
\tableofcontents

%%% List of figures in the thesis
\openright
\listoffigures

%%% List of tables in the thesis (optionally)
\clearpage
\listoftables

%%% List of abbreviations and glossary in the thesis (optionally)
\chapter*{List of abbreviations}

\begin{multicols}{2}
    \raggedright
    \begin{description}
        \item [\namedlabel{itm:html}{HTML}] Hypertext Markup Language
        \item [\namedlabel{itm:ux}{UX}] User eXperience
        \item [\namedlabel{itm:ui}{UI}] User Interface
        \item [\namedlabel{itm:gui}{GUI}] Graphical User Interface
        \item [\namedlabel{itm:rtf}{RTF}] Rich Text Format
        \item [\namedlabel{itm:xml}{XML}] eXtended Markup Language
        \item [\namedlabel{itm:cli}{CLI}] Command Line Interface
        \item [\namedlabel{itm:api}{API}] Application Programming Interface
    \end{description}
\end{multicols}


\chapter*{Glossary}

\begin{multicols}{2}
    \raggedright
    \begin{description}
        \item [\namedlabel{gloss:dotnetlabel}{.NET}] A free, cross-platform, open source developer platform for building many different types of applications
        \item [\namedlabel{gloss:aspnetcore}{ASP.NET Core}] A cross-platform, high-performance, open-source framework for building modern, cloud-enabled, Internet-connected apps
        \item [\namedlabel{gloss:nuget}{NuGet}] Package manager for .NET projects
    \end{description}
\end{multicols}

\pagestyle{fancy}
%%% The individual chapters of the thesis are stored in separate files for clarity
\chapter*{Introduction}
\addcontentsline{toc}{chapter}{Introduction}

The motivation for creating a custom tool for generating documentation from .NET libraries is a personal need to provide easy access to code documentation. Since most open source projects are hosted either on GitHub.com\footnote{A Microsoft-owned Git hosting platform} or GitLab.com\footnote{An independent Git hosting platform}, it makes sense to utilize said platforms for hosting source code documentation for consistency.

When attempting to find tools that generate documentation, none had the desired extensibility and were mainly limited to only creating static \ref{itm:html} pages. Since Git hosting platforms predominantly\footnote{GitLab also supports RDoc, AsciiDoc, and Org} utilize Markdown for displaying formatted text, static HTML pages aren't directly supported.

And so, the idea to develop a custom tool that would create Markdown documentation from .NET libraries for Git platforms came to fruition. However, focusing only on one output format would be a waste of time and effort, as few to none would need such a tool. Thus, the result of the development should be a generic tool that allows anyone to modify it to output to any desirable format.

Accomplishing this task will require analyzing the currently available documentation tools, getting user feedback, planning out the project structure, and finally developing the project itself.

\chapter{How .NET documentation works}

Source files of a language based on the \ref{gloss:dotnetlabel} platform can have structured comments that produce documentation for the types defined in those source files. The respective language compiler\footnote{C\# and VB are compiled by Roslyn, while F\# has its own specific F\# compiler} produces an \ref{itm:xml} file that contains structured data representing said comments and the type signatures.
\cite{wagner_xml_2022}

Documentation comments start with \lstinline[language=csh]{///} followed by \ref{itm:xml} tags. In C\#, for example, documenting a \lstinline[language=csh]{class} and its method would look like this:
\begin{lstlisting}[language=csh]
    /// <summary>
    /// This is a brief description of the class
    /// </summary>
    public class MyClass
    {
        /// <summary>
        /// This method executes some stuff
        /// </summary>
        public void Foo()
        {
            Bar();
        }
    }
\end{lstlisting}

\section{Documentation tags}
\subsection{General tags}
\subsection{Document members}
\subsection{Formatting tags}
\subsection{Documentation reuse tags}
\subsection{Reference tags}
\subsection{Generic types tags}

\chapter{Documentation tools}
\section{What is available?}
The current documentation tool market provides the following solutions:

\begin{enumerate}
    \item \href{https://github.com/dotnet/docfx}{DocFX}
    \item \href{https://www.doxygen.nl/}{Doxygen}
    \item \href{https://github.com/KirillOsenkov/SourceBrowser}{Source Browser}
    \item \href{https://github.com/lijunle/Vsxmd}{Vsxmd}
    \item \href{https://github.com/discosultan/vsdoc-2-md}{vsdoc-2-md}
\end{enumerate}

What follows are evaluations of each tool. This helped get a better understanding of the current offering and gain prospect on features and improvements that the custom tool can take advantage of.

\subsection{DocFX}

\textit{DocFX} is an open-source documentation generation tool developed by Microsoft that not only supports .NET languages C\#, F\#, and Visual Basic, but also Java, JavaScript, TypeScript, Python, and REST. Additionally, it can use raw Markdown files as input.

However, the \textit{DocFX} only outputs static \ref{itm:html} pages. The only available customizability is via templates for said static pages.

The tool is capable and is the go-to solution for \ref{gloss:dotnetlabel} projects, but it doesn't solve the task of outputting Markdown files.

\subsection{Doxygen}

\textit{Doxygen} is the industry-standard documentation generating tool originally made for C++ source code documentation. Nevertheless, it has over time added support for many popular programming languages such as C, C\#, Java, Python, and many more.

The tool provides an extensive set of supported output formats:
\begin{itemize}
    \item Static \ref{itm:html}
    \item \LaTeX\footnote{Document preparation system}
    \item Man pages\footnote{User manual type that is part of Unix operating systems \cite{credocs_limited_latex_2022}}
    \item \ref{itm:rtf}
    \item \ref{itm:xml}
\end{itemize}

\subsection{Source Browser}

\textit{Source Browser} is a tool that generates a website for browsing source code and its documentation. The tool is used by Microsoft, for example, to allow developers to browse the source code of \ref{gloss:dotnetlabel}.

The generated output is not fully static and has to be hosted on an \ref{gloss:aspnetcore} website to support searching.

\subsection{Vsxmd} \label{ssec:vsxmd}

\textit{Vsxmd} generates a single Markdown file for all types in a given assembly. Moreover, the tool has no \ref{itm:ui} and works as a \ref{gloss:nuget} package, that is added to the project designated for documentation generation. Thus, configuring this tool is done via \ref{gloss:dotnetlabel} project settings.

It is not possible to navigate the documentation, as no links are generated.

\subsection{Vsdoc-2-md}

\textit{Vsdoc-2-md} is an entirely unusable tool, as it is purely web-based and generates documentation from the provided \ref{itm:xml} documentation source file. The tool is limited to processing only one file at a time.

Just like \textit{\nameref{ssec:vsxmd}}, no links are generated; thus, it is not possible to navigate the documentation.

\section{What do users want?}
Surveying of potential users of the tool is necessary to ensure that the developed product is utilized by more than its developer.
Thus, creating a questionnaire was the first step to success. Its purpose was to figure out the following key points:
\begin{itemize}
    \item Whether users already use documentation generation tools, and if so, which ones
    \item What do users feel is missing from the tools they use
    \item What do users believe should be carried over to the new tool from the ones they use
    \item Whether users would integrate such tools into their \ref{itm:cicd} pipelines
    \item What output formats do users want
    \item What operating systems do users want the tool to work on
    \item Would users find a \ref{itm:gui} beneficial for such a tool
    \item Would users appreciate the tool being extensible via plugins
    \item Would users be interested in contributing to the project
\end{itemize}

\subsection{Questionnaire results}

What follows are the responses of the 22 participants of the questionnaire.

\subsubsection*{What output formats do you wish to have for your documentation?}

\begin{figure}[H]
    \centering
    \caption{What output formats do you wish to have for your documentation?}
    \begin{tikzpicture}
        \begin{axis}[
            symbolic x coords={Static HTML, PDF, Markdown, XML, LaTeX, Other},
            xtick=data,
            width=12cm,
            height=8cm,
            ymin=0,
            ymax=18,
            bar width=20pt,
            ylabel={Number of votes},
            enlarge y limits={value=0.2,upper},
            legend pos=north west,
            nodes near coords
        ]
            \addplot+[ybar] coordinates {
                (Static HTML,   16)
                (PDF,            9)
                (Markdown,       6)
                (XML,            7)
                (LaTeX,         1)
                (Other,          1)
            };
        \end{axis}
    \end{tikzpicture}
\end{figure}

\textit{Other responses:}
\begin{itemize}
    \item AsciiDoc
\end{itemize}

\subsubsection*{Would you integrate such a tool in your CI/CD process?}

\begin{figure}[H]
    \centering
    \caption{Would you integrate such a tool in your CI/CD process?}
    \begin{tikzpicture}
        \pie[sum=auto , after number=]{9/Yes, 2/No, 11/Maybe}
    \end{tikzpicture}
\end{figure}

\subsubsection*{How crucial is it for the tool to support these platforms?}

\begin{figure}[H]
    \centering
    \caption{How crucial is it for the tool to support these platforms?}
    \begin{tikzpicture}
        \begin{axis}[
            ybar stacked,
            bar width=20pt,
            height=10cm,
            width=8cm,
            nodes near coords,
            enlargelimits=0.15,
            legend style={at={(0.5,-0.25)},
            anchor=north,legend columns=-1},
            ylabel={Number of participants},
            symbolic x coords={Windows, Linux, MacOS},
            xtick=data,
            x tick label style={rotate=45,anchor=east},
            ]
        \addplot+[ybar] plot coordinates {(Windows,17) (Linux,4) (MacOS,1)};
        \addplot+[ybar] plot coordinates {(Windows,4) (Linux,16) (MacOS,2)};
        \addplot+[ybar] plot coordinates {(Windows,1) (Linux,2) (MacOS,19)};
        \legend{\strut First choice, \strut Middle choice, \strut Last choice}
        \end{axis}
    \end{tikzpicture}
\end{figure}

\subsubsection*{Would it be beneficial to have a GUI instead of only a console?}

\begin{figure}[H]
    \centering
    \caption{Would it be beneficial to have a GUI instead of only a console?}
    \begin{tikzpicture}
        \pie[sum=auto , after number=]{12/Yes, 9/No, 1/Other}
    \end{tikzpicture}
\end{figure}

\textit{Other responses:}
\begin{itemize}
    \item I don't want separate docs, readable code is better
\end{itemize}

\subsubsection*{Would you appreciate the tool being extensible via plugins?}

\begin{figure}[H]
    \centering
    \caption{Would you appreciate the tool being extensible via plugins?}
    \begin{tikzpicture}
        \pie[sum=auto , after number=]{17/Yes, 5/No}
    \end{tikzpicture}
\end{figure}

\subsubsection*{Are you using tools for documentation generation?}

\begin{figure}[H]
    \centering
    \caption{Are you using tools for documentation generation?}
    \begin{tikzpicture}
        \pie[sum=auto , after number=]{8/Yes, 14/No}
    \end{tikzpicture}
\end{figure}

\subsubsection*{What documentation generating tools do you use?}

\begin{figure}[H]
    \centering
    \caption{What documentation generating tools do you use?}
    \begin{tikzpicture}
        \pie[sum=auto , after number=]{3/DocFX, 2/Doxygen, 1/Slate, 1/Hugo}
    \end{tikzpicture}
\end{figure}

\subsubsection*{What features from those tools you'd like to see in the resulting project?}

\textit{Responses:}
\begin{itemize}
    \item \ref{itm:gui} for less technical people
    \item Links to external \ref{gloss:dotnetlabel} type references
    \item Easy templating
    \item Extended Markdown - custom shortcodes
    \item Fast build speed for \ref{itm:cicd}
    \item A capable \ref{itm:cli}
    \item Class structures
\end{itemize}

\subsubsection*{What features are missing in your current tool and would make you consider using the resulting project, if they were to be provided by it?}

\textit{Responses:}
\begin{itemize}
    \item Better support for diagrams (like mermaid.js)
    \item A tool that could build a static site with a) editorial (Markdown) content b) \ref{gloss:dotnetlabel} \ref{itm:api} reference c) RESTful API documentation from \ref{gloss:aspnetcore} projects and/or OpenAPI definitions all under one umbrella
\end{itemize}

\subsubsection*{Would you be interested in contributing to this project?}

\begin{figure}[H]
    \centering
    \caption{Would you be interested in contributing to this project?}
    \begin{tikzpicture}
        \pie[sum=auto , after number=]{14/No, 6/Maybe}
    \end{tikzpicture}
\end{figure}

\subsection{Questionnaire evaluation}

\chapter{Development strategy}
\chapter{M1: Proof of concept}

Before investing an unknown amount of time, it is best to verify whether the desired result is achievable. Developing a proof of concept application will give insight into the complexity of the task and how to tackle it.

The most crucial goal of this milestone is to find an optimal project structure. Neglecting this important aspect will introduce code quality degradation as the project grows.

For more see \ref{secProofOfConcept}.

\section{Prerequisites}

https://gitlab.com/hailstorm75/markdoc

\section{Development}

\section{Testing}

\section{Evaluation}
% \include{...}
% \include{...}
\chapter*{Summary}
\addcontentsline{toc}{chapter}{Summary}

Závěr je povinnou částí bakalářské/diplomové práce. Obsahuje shrnutí práce a vyjadřuje se k míře splnění cíle, který byl v práci stanoven, případně shrnuje odpovědi na otázky, které byly položeny v úvodu práce. 

Závěr k diplomové práci musí být propracovanější -- podrobněji to je uvedeno v Náležitostech diplomové práce v rámci Intranetu pro studenty FIS.

Závěr je vnímán jako kapitola (chapter), která začíná na samostatné stránce a která má název Závěr.  Název Závěr se nečísluje. Samotný text závěru je členěn do odstavců.

%%% Bibliography
%% Toto platí v případě použití samostatné bibliografické databáze
\printbibliography[title={\bibname},heading={bibintoc}]

%% Toto platí v případě použití prostředí thebibliography
%% Pro sestavení citačních údajů lze doporučit:
%%     https://knihovna.vse.cz/citace/priklady/
%%     https://www.citace.com/
%\openright
%\phantomsection
%\addcontentsline{toc}{chapter}{\bibname}
%\begin{thebibliography}{99}
%\bibitem{Cermak2018}ČERMÁK, Radim, SMUTNÝ, Zdeněk. A Framework for Cultural Localization of Websites and for Improving Their Commercial Utilization. In:  \emph{Global Observations of the Influence of Culture on Consumer Buying Behavior} [online]. Hershey~: IGI Global, 2018, s. 206--232. ISBN 978-1-5225-2727-5. DOI: 10.4018/978-1-5225-2727-5.
%
%\bibitem{Hladik2018}HLADÍK, Milan, ČERNÝ, Michal. The Shape of the Optimal Value of a Fuzzy Linear Programming Problem. In: \emph{Fuzzy Logic in Intelligent System Design} [online]. Cancum, 16.10.2017 -- 18.10.2017. Cham~: Springer, 2018, s. 281--286. Advances in Intelligent Systems and Computing 648. ISBN 978-3-319-67136-9. DOI: 10.1007/978-3-319-67137-6\_31.
%
%\bibitem{Jasek2018}JAŠEK, Pavel, VRANÁ, Lenka, ŠPERKOVÁ, Lucie, SMUTNÝ, Zdeněk, KOBULSKÝ, Marek. Modeling and Application of Customer Lifetime Value in Online Retail. \emph{Informatics} [online]. 2018, roč. 5, č. 1. 22 s. eISSN 2227-9709. DOI: 10.3390/informatics5010002. Dostupné také z: \url{http://www.mdpi.com/2227-9709/5/1/2/pdf}.
%
%\bibitem{Pecakova2018}PECÁKOVÁ, Iva. \emph{Statistika v terénních průzkumech}. 3. přeprac. vyd. Praha~: Professional Publishing, 2018. 254 s. ISBN 978-80-88260-10-3.
%\end{thebibliography}


%%% Attachments to thesis, if any. Each attachment must be referenced at
%%% least once in your own text. The appendices are numbered.
\part*{\Prilohy}
\appendix
\chapter{Formulář v plném znění}


\chapter{Zdrojové kódy výpočetních procedur}


% \include{...}
% \include{...}

\end{document}
